\documentclass[11pt]{article}
\usepackage[margin=1in]{geometry}
\usepackage[T1]{fontenc}
\usepackage{amsmath,amssymb,amsthm,mathtools}
\usepackage{mathrsfs}
\usepackage{microtype}
\usepackage{enumitem}
\usepackage{xcolor}
\usepackage[colorlinks=true,linkcolor=blue!55!black,
  citecolor=blue!55!black,urlcolor=blue!55!black]{hyperref}
\numberwithin{equation}{section}
\newtheorem{theorem}{Theorem}[section]
\newtheorem{lemma}[theorem]{Lemma}
\newtheorem{proposition}[theorem]{Proposition}
\newtheorem{corollary}[theorem]{Corollary}
\theoremstyle{definition}
\newtheorem{definition}[theorem]{Definition}
\newtheorem{remark}[theorem]{Remark}
\newtheorem{example}[theorem]{Example}
\newcommand{\cA}{\mathcal A}
\newcommand{\cD}{\mathcal D}
\newcommand{\cO}{\mathcal O}
\newcommand{\cQ}{\mathcal Q}
\newcommand{\cT}{\mathcal T}
\newcommand{\ex}{\operatorname{ex}}

\title{A proof of Kalai's tight-tree conjecture}
\author{Louis DeBiasio}
\date{September 6, 2026}

\hypersetup{pdfauthor={Louis DeBiasio}, pdftitle={A proof of Kalai's tight-tree conjecture}}

\begin{document}
\maketitle

\begin{abstract}
We prove Kalai's conjecture for tight uniform hypergraph trees.
If $T$ is a tight $r$-uniform tree with $m$ edges, where $r\geq2$,
then every $T$-free $r$-uniform hypergraph $H$ satisfies
$r e(H)\leq(m-1)|\partial H|$, where $\partial H$ is the family
of $(r-1)$-sets contained in its edges.  The proof extends the
permutation-counting argument of the recent Erd\H{o}s--S\'os proof
obtained in the FrontierMath Erd\H{o}s benchmark.  We also give a direct proof for
tight paths and a Lean~4 formalization of the general theorem.
\end{abstract}

\section*{Acknowledgments}
ChatGPT~6 Astra was used to develop the argument, write the exposition,
and produce the accompanying Lean formalization.  The author's role
was minimal, consisting mostly of asking questions towards trying to understand
the recent Erd\H{o}s--S\'os proof and its possible extensions.

\section{Introduction}
\label{sec:introduction}

In 1984, Gil Kalai \cite{Kalai1984} proposed a hypergraph extension
of the Erd\H{o}s--S\'os conjecture.  He conjectured that every
$n$-vertex $r$-uniform hypergraph with more than
$\frac{m-1}{r}\binom n{r-1}$ edges contains every tight $r$-uniform
tree with $m$ edges.  The conjecture was published by Frankl and
F\"uredi \cite[Conjecture~3.6]{FF1987}, who attribute it to a 1984
communication from Kalai.  We prove the conjecture, in the sharper
form that replaces $\binom n{r-1}$ by the number of $(r-1)$-sets
contained in edges of the hypergraph.

Our proof is based on the new proof of the Erd\H{o}s--S\'os conjecture
obtained in the FrontierMath Erd\H{o}s benchmark and recorded in
\cite{Erdos548Repository}.\footnote{The repository is maintained by
Tom Adamczewski.  Its provenance account credits the graph proof to
GPT-6 Astra, working autonomously in the FrontierMath Erd\H{o}s
benchmark of Adamczewski and Bloom.}
We adapt that proof's counting of marked vertex permutations and
induction on the target tree; we do not claim an independent origin
for this method.  In the graph argument, copies are rooted at a
specified vertex.  For the hypergraph argument, we prescribe the
images of all vertices in one edge.  This allows copies of branches
attached through different faces to agree on their common edge.
The counting inequalities and the required decomposition are proved
here for hypergraphs, rather than deduced by applying the graph
theorem as a black box.

\subsection{Definitions and main result}

All hypergraphs are finite and simple.  An $r$-uniform hypergraph, or
\emph{$r$-graph}, has edges that are $r$-element vertex sets.  A copy is
injective and edge-preserving, but need not be induced.

\begin{definition}
A \emph{tight $r$-tree} with $m\geq1$ edges admits an edge ordering
$e_1,\ldots,e_m$ such that, for every $j\geq2$, there is a vertex
$x_j\in e_j\setminus\bigcup_{h<j}e_h$ and an index $p(j)<j$ with
$e_j\setminus\{x_j\}\subseteq e_{p(j)}$.
The vertex set of a tight tree is the union of its edges; in particular,
there are no additional isolated vertices.  Such a tree has $m+r-1$
vertices.
The \emph{tight path} $P_m^{(r)}$ has distinct vertices
$x_1,\ldots,x_{m+r-1}$ and edges
\[
 \{x_j,x_{j+1},\ldots,x_{j+r-1}\},\qquad 1\leq j\leq m.
\]
\end{definition}

Throughout the paper, $m$ denotes the number of \emph{edges} of the
target.  We write $e(H)=|E(H)|$.

The $(r-1)$-\emph{shadow} of an $r$-graph $H$ is
\[
 \partial H=\{F: |F|=r-1,\ F\subseteq e\text{ for some }e\in E(H)\}.
\]
Our main result is the following.

\begin{theorem}[Tight-tree shadow bound]
\label{thm:main}
Let $r\geq2$ and $m\geq1$, and let $T$ be a tight $r$-tree with
$m$ edges.  Every $T$-free $r$-graph $H$ satisfies
\begin{equation}
 r e(H)\leq(m-1)|\partial H|.
 \label{eq:main-shadow}
\end{equation}
\end{theorem}

Since $|\partial H|\leq\binom{|V(H)|}{r-1}$, this proves Kalai's
conjecture in its usual form; see \cite{FJKMV2019,Stein2024}.

\begin{corollary}[Kalai's bound]
\label{cor:kalai}
Let $r\geq2$ and $m\geq1$, and let $T$ be a tight $r$-tree with
$m$ edges.  Every $n$-vertex $T$-free $r$-graph $H$ satisfies
\[
 e(H)\leq\frac{m-1}{r}\binom{n}{r-1}.
\]
\end{corollary}

The shadow formulation appears as Conjecture~2.1 in
\cite{FJKMV2019}.  Its coefficient is sharp: for $m\geq2$, the complete
$r$-graph on $m+r-2$ vertices has no copy of $T$ and attains equality
in \eqref{eq:main-shadow}.  The case $r=2$ gives the corresponding
tree bound with the number of nonisolated host vertices in place of
$|V(H)|$.

\subsection{Related work}

Frankl and F\"uredi \cite[Theorem~3.8]{FF1987} proved Kalai's
conjecture for star-shaped tight trees: those containing an edge that
meets every other edge in $r-1$ vertices.
F\"uredi, Jiang, Kostochka, Mubayi, and Verstra\"ete
\cite{FJKMV2019} proved Kalai's bound for tight trees with at most four
edges and an asymptotic form for tight trees of bounded trunk size.
Their work also develops bounds in terms of the shadow of the host.
For tight paths, their geometric method in \cite{FJKMV2020} gives
improved upper bounds in arbitrary uniformity.  In particular,
\cite[Theorem~1.2]{FJKMV2020} gives
\[
 \ex(n,P_m^{(3)})\leq
 \frac12\left(m+\left\lfloor\frac{m-1}{3}\right\rfloor\right)
 \binom n2.
\]
Here $\ex(n,F)$ denotes the maximum number of edges in an $n$-vertex
$F$-free hypergraph of the same uniformity as $F$.
Stein \cite{Stein2024} proved Kalai's conjecture for $r$-partite
$r$-graphs.  Bright, Milans, and Porter \cite{BMP2025} study the related
problem of excluding tight paths under order-preserving embeddings.
No partiteness or order-preservation restriction is imposed here.

\subsection{Outline of the proof}

Consider a permutation $(v_1,\ldots,v_n)$ of the host vertices.
We mark position $i\geq r$ if $v_i$ together with the first $r-1$
vertices forms an edge.  We count pairs consisting
of a permutation and one of its marked positions.  For a fixed tree,
we count those pairs for which the first $i$ vertices contain a copy
of the tree, with one specified edge mapped to the marking edge.
The copy need not follow the order of the permutation, and each pair
is counted only once, even if it admits several copies.

The induction compares these counts for smaller trees.  To add a leaf,
we use a copy lying entirely before the current marked position;
the vertex at that position is then available for the new leaf.
This excludes at most one pair for each permutation.  To join two
trees along a common edge, we require their copies to agree on that
edge and use disjoint sets of vertices elsewhere.  An injection
obtained by reordering the permutation gives the needed counting
inequality.  Finally, we show that every tight tree can be decomposed
using these two operations.  The resulting inequality gives the
claimed bound when the host contains no copy of the tree.

Section~\ref{sec:three} gives the marked-position argument for
$3$-uniform paths, and Section~\ref{sec:paths-r} treats arbitrary
uniformity.  Sections~\ref{sec:anchored} and
\ref{sec:anchored-transfers} count copies with a specified edge and
prove the two counting inequalities.  Section~\ref{sec:full-master}
completes the induction.
Section~\ref{sec:examples} records an example and sharpness, and
Section~\ref{sec:formalization} describes the accompanying formal proof.

\section{Marked positions for tight paths}
\label{sec:three}

We first prove the theorem for $3$-uniform paths.  In this case, the
injection rearranges just three entries of a permutation.  The main
point is to use a path that avoids the vertex at the current marked
position, so that this vertex can extend the path.

\subsection{Marked permutations}

Fix a $3$-graph $H$ on $n\geq3$ vertices.  For distinct vertices $a,b$,
write
\[
 d_H(ab)=|\{z:\{a,b,z\}\in E(H)\}|.
\]
For a permutation $\pi=(v_1,\ldots,v_n)$ and $3\leq i\leq n$, call
$(\pi,i)$ \emph{marked} if $\{v_1,v_2,v_i\}\in E(H)$.  Its prefix
is the vertex set $S_i(\pi)=\{v_1,\ldots,v_i\}$.
Let $\cO(H)$ be the set of marked states and $M=|\cO(H)|$.

For $k\geq1$, let $\cA_k$ be the set of marked states whose prefixes
contain a tight $k$-edge path with vertex sequence
\begin{equation}
 v_1,v_2,y_3,\ldots,y_{k+2}.
 \label{eq:rooted-path}
\end{equation}
The starting pair is ordered, but the remaining path vertices need not
occur in the permutation in path order.  Support is an existence
condition inside the prefix vertex set.  Each supporting state is
counted once, independently of the number of witnesses.  Write
$A_k=|\cA_k|$.

\begin{lemma}[Basic counts]
\label{lem:normalization}
With the notation above,
\begin{equation}
 A_1=M=6e(H)(n-2)!.
 \label{eq:normalization}
\end{equation}
Let $\cQ$ be the set of permutations whose first pair belongs to
$\partial H$.  Then
\begin{equation}
 Q:=|\cQ|=2|\partial H|(n-2)!.
 \label{eq:Q}
\end{equation}
\end{lemma}

\begin{proof}
Every marked state supports a one-edge path starting at its first pair,
so $\cA_1=\cO(H)$.  For each fixed ordered pair $(a,b)$, there are
$(n-2)!$ permutations starting with $(a,b)$, and each contributes
$d_H(ab)$ marks.  Each edge contributes to six ordered pairs, giving
\[
 M=(n-2)!\sum_{a\ne b}d_H(ab)=6e(H)(n-2)!.
\]
Each shadow pair has two orders and each such order can begin
$(n-2)!$ permutations.  This proves \eqref{eq:Q}.
\end{proof}

\subsection{Extending a path}

If an earlier marked position already admits the required path, that
path avoids the vertex at the current marked position.  We therefore
discard the first such position in each permutation and extend the
paths at all remaining positions.  This gives the following inequality.

\begin{lemma}[Path extension]
\label{lem:transfer}
For every $k\geq1$,
\begin{equation}
 A_k\leq A_{k+1}+Q.
 \label{eq:transfer}
\end{equation}
\end{lemma}

\begin{proof}
For each permutation $\pi$ contributing to $\cA_k$, let $i_0(\pi)$
be its first position in $\cA_k$, and put
\[
 \cD_k=\{(\pi,i_0(\pi)):\pi\text{ contributes to }\cA_k\}.
\]
Such a permutation has a marked position, so its first pair lies in
$\partial H$.  Hence $|\cD_k|\leq Q$.

We construct an injection $\Phi_k:\cA_k\setminus\cD_k\to\cA_{k+1}$.
Take a state $(\pi,i)$ in its domain and write
\[
 \pi=(a,b,W,z,R),\qquad
 W=(v_3,\ldots,v_{i-1}),\quad z=v_i,\quad
 R=(v_{i+1},\ldots,v_n).
\]
Since $i>i_0(\pi)$, the prefix $S_{i_0(\pi)}(\pi)$ contains a path
\[
 a,b,y_3,\ldots,y_{k+2}
\]
and does not contain $z$.  The marking triple $\{z,a,b\}$ extends
this to a path
\begin{equation}
 z,a,b,y_3,\ldots,y_{k+2}
 \label{eq:extended-path}
\end{equation}
with $k+1$ edges.  Indeed, its first edge is $\{z,a,b\}$ and its
remaining consecutive triples are precisely the old path edges.
All vertices remain distinct.

Define
\begin{equation}
 \Phi_k\bigl((a,b,W,z,R),i\bigr)
       =\bigl((z,a,W,b,R),i\bigr).
 \label{eq:cycle-map}
\end{equation}
The new prefix has the same vertex set.  Its first pair is $(z,a)$,
the starting pair of \eqref{eq:extended-path}, and its last vertex is
$b$.  Thus its marking triple is still $\{z,a,b\}$.  The image is
therefore in $\cA_{k+1}$.

The retained index $i$ identifies the displaced entry $b$.
For any image $(\rho,i)$, the unique possible source ordering is
\begin{equation}
 (\rho_2,\rho_i,\rho_3,\ldots,\rho_{i-1},\rho_1,
       \rho_{i+1},\ldots,\rho_n).
 \label{eq:inverse}
\end{equation}
This inverse cycle on positions $1,2,i$ recovers the source without
using the earlier index $i_0(\pi)$ or choosing a path witness.  Hence
$A_k-|\cD_k|\leq A_{k+1}$ proves \eqref{eq:transfer}.
\end{proof}

\begin{theorem}[Tight-path shadow bound]
\label{thm:three}
For every integer $m\geq1$ and every $P_m^{(3)}$-free $3$-graph $H$,
\begin{equation}
 3e(H)\leq(m-1)|\partial H|.
 \label{eq:shadow-bound}
\end{equation}
In particular, for an $n$-vertex host,
$e(H)\leq\frac{m-1}{3}\binom n2$.
\end{theorem}

\begin{proof}
If $n<3$ there are no edges, so assume $n\geq3$.
Iterating Lemma~\ref{lem:transfer} gives
\[
 6e(H)(n-2)!=A_1\leq A_m+(m-1)Q.
\]
The forbidden-path hypothesis gives $A_m=0$; substitution of
\eqref{eq:Q} and division by $2(n-2)!$ gives
\eqref{eq:shadow-bound}.  For $m=1$ the iteration is empty and
$A_1=0$, which gives the same conclusion.  Finally,
$|\partial H|\leq\binom n2$.
\end{proof}

\subsection{The role of the mark}

The path and the marking edge need not start with the same edge.
The path only has to lie among the first $i$ vertices and start with
the ordered pair $(v_1,v_2)$.  The mark supplies another edge through
that pair.  Moving $z$ to the first position and $b$ to the marked
position preserves both requirements.  The map does not depend on
which path is chosen; the earlier marked position ensures that a path
avoiding $z$ exists.

\begin{example}
For five distinct vertices $a,b,c,d,z$, write $abc$ for
$\{a,b,c\}$, and similarly for other edges.  Suppose $H$ contains
$abc,bcd,zab$, and take
\[
 \pi=(a,b,c,d,z),\qquad i=5.
\]
The prefix at position $4$ contains the two-edge path $a,b,c,d$.
However, position $4$ is marked only if $abd\in E(H)$.  Without that
edge, position $5$ is the first supporting mark and is discarded.
If $abd\in E(H)$, the map is defined at position $5$ and gives
\[
 (a,b,c,d,z)\longmapsto(z,a,c,d,b),
\]
whose full prefix supports $z,a,b,c,d$ and whose mark is $zab$.
Thus an earlier supporting prefix alone is not sufficient: the proof
uses an earlier supporting \emph{marked state}.
\end{example}

\begin{remark}[Boundary cases and equality]
An empty host has $A_k=M=Q=0$.  For $n=3$, every permutation has at
most one mark, so every supporting pair is discarded.
For the complete $3$-graph on $s=m+1\geq3$ vertices, there is no
$P_m^{(3)}$, and
\[
 3\binom{s}{3}=(m-1)\binom{s}{2}.
\]
Thus the coefficient in the shadow bound cannot be reduced.  More
precisely, every permutation has $s-2$ marks and exactly $s-k-1$
of them support a $k$-edge path, for $1\leq k\leq s-1$.  Hence
$A_k=(s-k-1)s!$ and $A_k-A_{k+1}=Q=s!$ throughout the nontrivial
iteration.  Adding isolated vertices does not change the shadow or edge
count and only changes the factorial normalizations.
\end{remark}

\section{Tight paths in arbitrary uniformity}
\label{sec:paths-r}

The preceding proof extends by replacing the three-cycle with an
$r$-cycle of entries.  We give the map explicitly.
Fix $r\geq2$ and $n\geq r$.  Mark $(\pi,i)$, $i\geq r$, when
\[
 \{v_1,\ldots,v_{r-1},v_i\}\in E(H).
\]
Let $A_k^{(r)}$ count the marked states whose prefixes support a tight
$k$-edge path starting with the ordered tuple
$(v_1,\ldots,v_{r-1})$.  The same normalization gives
\begin{equation}
 A_1^{(r)}=r!\,e(H)(n-r+1)!,\qquad
 Q_r=(r-1)!\,|\partial H|(n-r+1)!.
 \label{eq:r-normalization}
\end{equation}
Here $Q_r$ counts permutations whose first $(r-1)$-set belongs to the
shadow.

Discard the first supporting marked state of each permutation.  At a
remaining state, write
\[
 \pi=(a_1,\ldots,a_{r-1},W,z,R).
\]
An earlier path beginning $a_1,\ldots,a_{r-1}$ avoids $z$, and
prepending $z$ adds exactly the marking edge.  Apply
\begin{equation}
 \bigl((a_1,\ldots,a_{r-1},W,z,R),i\bigr)
 \longmapsto
 \bigl((z,a_1,\ldots,a_{r-2},W,a_{r-1},R),i\bigr).
 \label{eq:r-cycle}
\end{equation}
The prefix set and marking $r$-set are unchanged.  The new first
$(r-1)$-tuple is the starting tuple of the extended path.
The inverse takes the entries at positions $2,\ldots,r-1,i$ as the
old root tuple, restores the entry at position $1$ to position $i$,
and leaves all other entries fixed.  Thus
\[
 A_k^{(r)}\leq A_{k+1}^{(r)}+Q_r.
\]
When $r=2$, the tuple $(a_1,\ldots,a_{r-2})$ is empty and the map
simply exchanges the first entry and the mark.

\begin{corollary}
\label{cor:all-r}
If an $r$-graph $H$ contains no tight path with $m$ edges, then
\[
 r e(H)\leq(m-1)|\partial H|.
\]
Consequently it satisfies Corollary~\ref{cor:kalai}
for the target $P_m^{(r)}$.
\end{corollary}

\begin{proof}
For $n<r$ there are no edges.  Otherwise iterate the preceding
inequality, use $A_m^{(r)}=0$, and divide
\eqref{eq:r-normalization} by $(r-1)!(n-r+1)!$.
\end{proof}

\section{Copies with a specified edge}
\label{sec:anchored}

For paths, it was enough to specify the first $r-1$ vertices of a
copy.  A tight tree may have several edges attached through different
$(r-1)$-subsets of the same edge.  For example, the tight
$3$-tree
\[
 123,\qquad124,\qquad135,\qquad236
\]
has edges attached through all three pairs of $123$.  To combine
copies of these branches, we must make them agree on the images of
all three vertices of $123$, not just two of them.

We therefore choose an edge of the target, called its \emph{root edge},
and specify the image of each of its vertices.  Unlike the condition
used for paths, we now require this chosen edge to map to the marking
edge.  Copies of two branches can then be made to agree on their common
edge.  In the proof of Lemma~\ref{lem:anchored-glue}, we arrange for
their other vertices to lie in disjoint parts of the permutation.

\subsection{Definitions and basic counts}

Fix $r\geq2$ and an $r$-graph $H$ on $n\geq r$ vertices.  Let
\[
 \cO_r(H)=\{(\pi,i):\pi=(v_1,\ldots,v_n),\ r\leq i\leq n,
       \ \{v_1,\ldots,v_{r-1},v_i\}\in E(H)\}.
\]
Put $S_i(\pi)=\{v_1,\ldots,v_i\}$ and let $\cQ_r$ be the set of
permutations whose first $(r-1)$-set lies in $\partial H$.  Throughout
the general-tree argument write
\begin{equation}
 M=|\cO_r(H)|=r!\,e(H)(n-r+1)!,\qquad
 Q=|\cQ_r|=(r-1)!\,|\partial H|(n-r+1)!.
 \label{eq:anchored-normalization}
\end{equation}
To verify \eqref{eq:anchored-normalization}, each ordered
$(r-1)$-tuple begins
$(n-r+1)!$ permutations, each with one mark for every edge containing
that tuple.  Summing over the ordered tuples gives $r!e(H)$.  Counting
the orders of each shadow face gives the formula for $Q$.

\begin{definition}[Support with a specified edge]
Let $T$ be an $r$-graph with a distinguished ordered edge
$\mathbf e=(u_1,\ldots,u_r)$, whose underlying set is $e$.
A marked state $(\pi,i)$ \emph{supports $(T,\mathbf e)$} if there is
an injective edge-preserving map $\varphi:V(T)\to S_i(\pi)$ with
\begin{equation}
 \varphi(u_j)=v_j\quad(1\leq j<r),\qquad
 \varphi(u_r)=v_i.
 \label{eq:anchored-support}
\end{equation}
Let $\cA_H(T,\mathbf e)$ be the set of these states and put
$R_H(T,\mathbf e)=|\cA_H(T,\mathbf e)|$.  Each state is counted once,
regardless of the number of witnessing embeddings.  We omit the host
subscript when it is fixed.
\end{definition}

\begin{lemma}[Independence of the ordering of the root edge]
\label{lem:root-order}
The count $R(T,\mathbf e)$ depends on the root edge $e$, but not on
the ordering of its vertices.  We may therefore denote it by $R(T,e)$.
\end{lemma}

\begin{proof}
Fix the marked position $i$.  If the target root order is changed by a
permutation of its $r$ entries, apply the same permutation to the host
entries in positions $1,\ldots,r-1,i$.  All other entries and the index
$i$ remain fixed.  The prefix vertex set and the marking edge are
unchanged.  The same embedding sends each root vertex to its required
position after the change.  The inverse
permutation reverses the construction, so this is a bijection for each
$i$, and hence for their disjoint union.
\end{proof}

Thus reordering the root vertices does not change the total count.
It can change whether an individual pair $(\pi,i)$ supports the
target, so we will specify the root orders when comparing sets of pairs.

\section{Extending and joining copies}
\label{sec:anchored-transfers}

We need two inequalities: one for adding a leaf edge, and one for
joining two hypergraphs along an edge.  The second proof is easier
to write after moving the marking edge to the first $r$ positions.
We first show that this rearrangement leaves the counts unchanged.

\subsection{A new leaf edge and a change of root}

\begin{lemma}[Adding a leaf edge]
\label{lem:anchored-leaf}
Let $S$ have a distinguished edge $f=F\cup\{p\}$, where
$|F|=r-1$.  Let $z\notin V(S)$, and form $T$ by adding $z$ and
the edge $e=F\cup\{z\}$ to $S$.  Then
\begin{equation}
 R(S,f)\leq R(T,e)+Q.
 \label{eq:anchored-leaf}
\end{equation}
Thus the distinguished edge changes from $f$ to the newly added edge $e$.
\end{lemma}

\begin{proof}
Order $F$ as $(u_1,\ldots,u_{r-1})$ and use the orders
$\mathbf f=(u_1,\ldots,u_{r-1},p)$ and
$\mathbf e=(u_1,\ldots,u_{r-1},z)$.  Lemma~\ref{lem:root-order}
allows these choices without changing either count.

For each permutation $\pi$ having a state in $\cA(S,\mathbf f)$,
discard its first such state, at position $i_0(\pi)$.  Let $\cD$ be
the set of discarded states.  Each contributing permutation has a mark,
so belongs to $\cQ_r$; hence $|\cD|\leq Q$.

Take $(\pi,i)\in\cA(S,\mathbf f)\setminus\cD$.  An earlier state
$(\pi,i_0)$ supports $S$, with $u_j\mapsto v_j$ for $j<r$ and
$p\mapsto v_{i_0}$.  This earlier embedding avoids $v_i$ because
$i_0<i$.  Extend it by setting $z\mapsto v_i$.  The only added target
edge maps to the current marking edge
$\{v_1,\ldots,v_{r-1},v_i\}$.  The extended embedding is injective
and sends each vertex of $e$ to its required position.

Consequently the \emph{identity map on states} gives the inclusion
\[
 \cA(S,\mathbf f)\setminus\cD\ \subseteq\ \cA(T,\mathbf e).
\]
Taking cardinalities proves \eqref{eq:anchored-leaf}.  The current
state's witness for $S$ need not be used; it is the earlier witness
that guarantees avoidance of the new vertex.
\end{proof}

\subsection{Moving the marking edge to the front}

Consider the following alternative set of pairs:
\[
 \mathcal B_r(H)=\{(\rho,i):\rho=(w_1,\ldots,w_n),\ r\leq i\leq n,
                         \ \{w_1,\ldots,w_r\}\in E(H)\}.
\]
The first $r$ entries form an ordered edge, and $i$ specifies how many
vertices may be used in a copy.  Every value $r,\ldots,n$ is allowed.
We indicate the end of the prefix by a vertical bar after position $i$,
which we call the \emph{divider}.  The bijection between the two
sets of pairs is
\begin{equation}
 \begin{split}
 \Theta:\cO_r(H)&\longrightarrow\mathcal B_r(H),\\
 ((a_1,\ldots,a_{r-1},W,c,C),i)
   &\longmapsto ((a_1,\ldots,a_{r-1},c,W,C),i),
 \end{split}
 \label{eq:edge-front}
\end{equation}
where $c$ was at position $i$ and $|W|=i-r$.  The index is retained.
For $(\rho,i)\in\mathcal B_r(H)$, the inverse ordering is
\begin{equation}
 (w_1,\ldots,w_{r-1},w_{r+1},\ldots,w_i,w_r,
                                      w_{i+1},\ldots,w_n).
 \label{eq:edge-front-inverse}
\end{equation}
Empty blocks are omitted, so the formula also applies when $i=r$.

The prefix set is unchanged.  Under $\Theta$, condition
\eqref{eq:anchored-support} becomes $u_j\mapsto w_j$ for every
$1\leq j\leq r$.  Let $\mathcal B(T,\mathbf e)$ denote the supporting
states in this representation.  Thus
\begin{equation}
 |\mathcal B_r(H)|=M,\qquad
 |\mathcal B(T,\mathbf e)|=R(T,e).
 \label{eq:front-counts}
\end{equation}
The first identity also follows directly:
there are $r!e(H)(n-r)!$ orderings with an edge first, and each has
$n-r+1$ possible values of $i$.

\subsection{Joining copies along a common edge}

\begin{lemma}[Joining along an edge]
\label{lem:anchored-glue}
Suppose $T=T_1\cup T_2$, the edge $e$ belongs to both $T_1$ and $T_2$,
and $V(T_1)\cap V(T_2)=e$.  Then
\begin{equation}
 R(T_1,e)+R(T_2,e)\leq M+R(T,e).
 \label{eq:anchored-glue}
\end{equation}
The inequality holds for arbitrary $r$-graphs $T_1,T_2$ with the stated
intersection; they need not be tight trees.
\end{lemma}

\begin{proof}
Fix the same ordering $\mathbf e$ in all three targets and work in
$\mathcal B_r(H)$.  We construct an injection
\begin{equation}
 \mathcal B(T_1,\mathbf e)\setminus\mathcal B(T,\mathbf e)
 \longrightarrow
 \mathcal B_r(H)\setminus\mathcal B(T_2,\mathbf e).
 \label{eq:anchored-glue-injection}
\end{equation}

Take a source state $(\rho,i)$, and let $a$ be the least index in
$\{r,\ldots,i\}$ at which its prefix supports $T_1$ with the prescribed
root-edge correspondence.  Write
\[
 \rho=(E_0,A,B,C),\qquad
 E_0=(w_1,\ldots,w_r),\quad
 A=(w_{r+1},\ldots,w_a),\quad B=(w_{a+1},\ldots,w_i).
\]
Here $E_0$ lists the root-edge images, and $A,B,C$ are consecutive
lists of the remaining vertices.  The original prefix ends after $B$.
Set $b=r+|B|$ and define
\begin{equation}
 ((E_0,A,B,C),i)\longmapsto ((E_0,B,A,C),b).
 \label{eq:anchored-rotation}
\end{equation}
In divider notation this is
\[
 (E_0,A,B\mid C)\longmapsto(E_0,B\mid A,C).
\]
The image is a valid state of $\mathcal B_r(H)$, including when $B$
is empty, since then $b=r$.

By the choice of $a$, the vertices in $E_0,A$ contain a copy of $T_1$
with the specified images of the root vertices.  If the image prefix
$E_0,B$ supported $T_2$, the
two embeddings would agree on every vertex of $e$ and be disjoint
outside $e$.  Indeed, the non-root vertices of the first copy lie in
$A$ and those of the second lie in $B$.  Their union would therefore
embed $T$ into the original prefix, contrary to the source condition.
This proves that the image is in the right-hand side of
\eqref{eq:anchored-glue-injection}.

To recover the source from an image, first read $E_0$ from the first
$r$ positions and $B$ from the rest of the image prefix.  Starting
immediately after the image divider, find the shortest initial block
$A'$ such that $E_0,A'$ supports $T_1$ with its prescribed root edge.
Allow $A'$ to be empty.  The actual block $A$ works.  No shorter
initial block works, since such a block would be a proper initial
segment of $A$ and would contradict the original minimality of $a$.
Thus $A'=A$.  This test uses $E_0$ and the block after the divider;
it does \emph{not} include $B$ in the support test.  The remaining
entries are $C$, and the original index is
$i=r+|A|+|B|$.  Both the ordering and index have been recovered uniquely.

Finally, $\mathcal B(T,\mathbf e)\subseteq\mathcal B(T_1,\mathbf e)$.
Taking cardinalities in the injection and using
\eqref{eq:front-counts} gives
\[
 R(T_1,e)-R(T,e)\leq M-R(T_2,e),
\]
which is \eqref{eq:anchored-glue}.
\end{proof}

\begin{remark}[The case of an empty middle list]
If $B$ is empty, the image has $i=r$: its prefix consists just of the
root edge.  This value of $i$ is included in $\mathcal B_r(H)$, so
the map is still defined and its inverse still recovers $A$.
No pairs need to be discarded in this case.
\end{remark}

\section{Completing the induction}
\label{sec:full-master}

We now show that every tight tree can be decomposed into the two cases
treated above, whatever edge is chosen as the root.  We do this by
constructing an ordinary tree whose nodes are the hyperedges.  This
auxiliary tree records how the hyperedges are attached to one another.

\subsection{A tree on the hyperedges}

\begin{lemma}[An auxiliary tree on the edges]
\label{lem:edge-tree}
Let $T$ be a tight $r$-tree.  There is a tree $\mathscr T$ whose nodes
are the hyperedges of $T$ such that:
\begin{enumerate}[label=\textup{(\roman*)},leftmargin=*]
\item adjacent nodes correspond to edges meeting in exactly $r-1$ vertices;
\item for each vertex $x$ of $T$, the nodes whose edges contain $x$
      form a connected subtree of $\mathscr T$;
\item after rooting $\mathscr T$ at any node, every parent-before-child
      ordering is a tight-tree construction order for $T$.
\end{enumerate}
The edges indexed by any nonempty connected subtree of $\mathscr T$,
with their union as vertex set, also form a tight $r$-tree.
\end{lemma}

\begin{proof}
Choose a tight-tree construction order $e_1,\ldots,e_m$ and one
permitted parent $e_{p(j)}$ for each $j\geq2$.  Join $e_j$ to
$e_{p(j)}$.  Each new node attaches to a preceding node by one edge,
so these connections form a tree.  The child introduces exactly one
new vertex, and all its other $r-1$ vertices belong to the parent;
this proves (i).

To prove (ii), follow the construction.  A vertex first occurs in a
single node.  Whenever a subsequent edge contains that existing vertex,
its parent must contain it too, since every child vertex other than the
newly introduced one belongs to the parent.  The nodes containing that
vertex therefore remain connected at every step.

Now root $\mathscr T$ at any node and order its nodes with parents
before children.  Consider a child edge $g$ with parent $f$, and let
$x$ be the unique vertex of $g\setminus f$.  If $x$ occurred in an
earlier edge $h$, then $h$ could not be a descendant of $g$, because
descendants occur after $g$.  The path from $g$ to $h$ would pass
through $f$.  Property (ii) would force $x\in f$, a contradiction.
Thus $x$ is new, and $g\setminus\{x\}\subseteq f$ proves (iii).

For a connected subtree, the nodes containing any given vertex are
again connected whenever nonempty: the unique path between two such
nodes lies in both the original vertex-subtree and the chosen subtree.
The same parent-before-child argument proves the final assertion.
\end{proof}

\begin{lemma}[Decomposition at a distinguished edge]
\label{lem:edge-decomposition}
Let $T$ have $m\geq2$ edges and distinguished edge $e$, and root
$\mathscr T$ from Lemma~\ref{lem:edge-tree} at $e$.
\begin{enumerate}[label=\textup{(\roman*)},leftmargin=*]
\item If $e$ has one neighbour $f$ in $\mathscr T$, the vertex of
      $e\setminus f$ occurs in no other edge of $T$.  Removing $e$
      and that vertex leaves a tight tree $S$ with $m-1$ edges.
\item If $e$ has at least two neighbours, $T$ can be written
      $T_1\cup T_2$, with both $T_j$ smaller tight trees containing
      $e$, such that
      \[
       V(T_1)\cap V(T_2)=e,\qquad e(T_1)+e(T_2)=m+1.
      \]
\end{enumerate}
\end{lemma}

\begin{proof}
In (i), if the vertex $x\in e\setminus f$ occurred in another node,
the path to that node would pass through $f$, contradicting
property (ii) of Lemma~\ref{lem:edge-tree}.  The remaining nodes form
a connected subtree, so they
give a tight tree by Lemma~\ref{lem:edge-tree}.

In (ii), partition the components of $\mathscr T-e$ into two nonempty
groups and restore $e$ to each group.  Each group with $e$ forms a
connected subtree, and hence defines a tight tree $T_j$.  If a vertex
occurs in both groups, property (ii) of Lemma~\ref{lem:edge-tree},
applied to the path through $e$, forces that vertex to belong to $e$.
Conversely, both trees contain
all vertices of $e$.  Their edge sets overlap in exactly the single
edge $e$, which gives the displayed edge-count identity.  Nonemptiness
of both groups ensures that each $T_j$ has fewer than $m$ edges.
\end{proof}

\subsection{The counting inequality}

\begin{theorem}[The main counting inequality]
\label{thm:anchored-master}
Let $r\geq2$, let $T$ be a tight $r$-tree with $m\geq1$ edges,
and distinguish any edge $e$ of $T$.  For every $r$-graph $H$ on
$n\geq r$ vertices,
\begin{equation}
 M\leq R_H(T,e)+(m-1)Q,
 \label{eq:anchored-master}
\end{equation}
where $M$ and $Q$ are given in \eqref{eq:anchored-normalization}.
\end{theorem}

\begin{proof}
Induct on $m$, simultaneously for all tight trees and all choices
of their distinguished edges.  If $m=1$, every marked state supplies
the prescribed copy of the sole edge, so $R_H(T,e)=M$.

Suppose $m\geq2$ and use the decomposition of
Lemma~\ref{lem:edge-decomposition}.  If $e$ has one neighbour $f$,
write $S$ for the remaining tight tree.  The edge $e$ is obtained
from $f$ by retaining its intersection with $e$ and adjoining the
vertex of $e$ absent from $S$.  Induction rooted at $f$, followed by
Lemma~\ref{lem:anchored-leaf}, gives
\[
 M\leq R_H(S,f)+(m-2)Q
   \leq R_H(T,e)+(m-1)Q.
\]

Otherwise write $T=T_1\cup T_2$ as in part (ii) and put
$m_j=e(T_j)$.  Add the induction inequalities rooted at $e$ and apply
Lemma~\ref{lem:anchored-glue}:
\[
 \begin{aligned}
 2M
 &\leq R_H(T_1,e)+R_H(T_2,e)+(m_1+m_2-2)Q\\
 &\leq M+R_H(T,e)+(m-1)Q.
 \end{aligned}
\]
Subtracting $M$ proves \eqref{eq:anchored-master}.  Root-order choices
made in the leaf step do not change the counts, by
Lemma~\ref{lem:root-order}.
\end{proof}

\begin{proof}[Proof of Theorem~\ref{thm:main}]
Write $n=|V(H)|$.  For $n<r$, the host has no edges and the inequality
is immediate.
Otherwise distinguish any edge $e$ of $T$.  The forbidden-copy
hypothesis gives $R_H(T,e)=0$.  Substitute
\eqref{eq:anchored-normalization} into
Theorem~\ref{thm:anchored-master} to obtain
\[
 r!\,e(H)(n-r+1)!
 \leq(m-1)(r-1)!\,|\partial H|(n-r+1)!.
\]
Division by the positive factorial factor $(r-1)!(n-r+1)!$ gives
\eqref{eq:main-shadow}.  This includes the one-edge target and the empty
host: the argument does not divide by $Q$.
\end{proof}

\section{An example and sharpness}
\label{sec:examples}

The following example illustrates why the copies must agree on every
vertex of their common edge.

\begin{example}[Branches through three different pairs]
\label{ex:three-faces}
Let $r=3$ and distinguish $e=123$ in the tight tree
\[
 E(T)=\{123,124,135,236\}.
\]
Write $T_1=\{123,124\}$, $T_2=\{123,135\}$, and
$T_3=\{123,236\}$, with the vertex sets understood to be their unions.
For each $j$, start with the single edge other than $123$ and apply
Lemma~\ref{lem:anchored-leaf}, adding $123$ as the new distinguished edge.
The new vertices are respectively $3,2,1$, so this gives
\[
 R_H(T_j,123)\geq M-Q\qquad(j=1,2,3).
\]
Two applications of Lemma~\ref{lem:anchored-glue} yield
\[
 R_H(T,123)\geq
 R_H(T_1,123)+R_H(T_2,123)+R_H(T_3,123)-2M
 \geq M-3Q.
\]
The copies agree on the images of $1,2,3$.  The injection in
Lemma~\ref{lem:anchored-glue} uses copies whose other vertices lie in
disjoint parts of the permutation, so their union is a copy of $T$.
The resulting
bound for this particular four-edge tree also follows from
\cite{FJKMV2019}.
\end{example}

\begin{proposition}[Sharpness of the shadow coefficient]
For every $r\geq2$, $m\geq2$, and tight $r$-tree $T$ with $m$
edges, there is a $T$-free $r$-graph attaining equality in
Theorem~\ref{thm:main}.
\end{proposition}

\begin{proof}
Put $s=m+r-2$ and take $H=K_s^{(r)}$.  Since $T$ has
$m+r-1=s+1$ vertices, it cannot embed in $H$.  Moreover,
\[
 r\binom{s}{r}=(s-r+1)\binom{s}{r-1}
             =(m-1)\binom{s}{r-1}.
\]
Here $|\partial H|=\binom{s}{r-1}$, giving the required equality.
\end{proof}

Disjoint unions of these complete hypergraphs also attain equality in
the shadow bound: a tight tree is connected, and the edge and shadow
counts add over the components.  Adding isolated host vertices leaves
both counts unchanged.

When $r=2$, the shadow consists of the singleton sets corresponding to
nonisolated vertices.  Theorem~\ref{thm:main} therefore reads
\[
 2e(H)\leq(m-1)|\{v\in V(H):d_H(v)>0\}|,
\]
and implies the Erd\H{o}s--S\'os bound.  As explained in the
introduction, the method is adapted from \cite{Erdos548Repository};
the graph theorem itself is not used as a black box.

\section{Lean formalization}
\label{sec:formalization}

The accompanying project in \texttt{lean/} proves the main theorem,
its binomial consequence, and the counting inequality in Lean~4.19.0.
It uses Lean's bundled \texttt{Std} library and does not depend on
Mathlib.  We describe the definitions and their relation to the
mathematical proof; instructions for checking the files are given in
\texttt{lean/README.md}.

\subsection{Definitions and theorem statements}

The structure \texttt{FiniteHypergraph} represents a finite simple
hypergraph by lists of vertices and edges.  The vertex list and each
edge list are duplicate-free, and no two stored edges are permutations
of one another.  Edge membership is defined up to permutation, so this
representation imposes no order-preservation condition on embeddings.
An \texttt{Embedding} is injective on the listed target vertices and
maps each target edge to a host edge; nonedges need not be preserved.

The definition \texttt{IsTightTree} is the edge-order definition from
Section~\ref{sec:introduction}: each edge after the first adds one
new vertex, and its other vertices lie in an earlier edge.  It also
requires uniformity and excludes isolated target vertices.  The
decomposition and counting inequalities are proved, not assumed.
The parameter $m$ is the number of edges, and the identity
$|V(T)|=m+r-1$ is also proved in Lean.

The principal declarations are
\begin{itemize}[leftmargin=*]
\item \texttt{Kalai.kalai\_master}, corresponding to
      Theorem~\ref{thm:anchored-master};
\item \texttt{Kalai.kalai\_shadow\_bound}, corresponding to
      Theorem~\ref{thm:main};
\item \texttt{Kalai.kalai\_bound}, corresponding to
      Corollary~\ref{cor:kalai}; and
\item \texttt{Kalai.kalai\_contains\_tree}, giving the equivalent
      embedding conclusion above the binomial threshold.
\end{itemize}
The inequalities use natural numbers, with the denominator $r$ cleared.
The code counts each shadow set once, regardless of how many edges
contain it.  It defines the binomial coefficient by Pascal's recurrence
and proves the corresponding subset-counting identity.  The statements
include empty hosts and hosts with fewer than $r$ vertices.

\subsection{Relation to the proof}

The formal proof establishes the two counting inequalities, the
injection for joining copies and its inverse, and the formulas for
$M$ and $Q$.  It also checks that reordering the chosen edge and
moving the marking edge to the front leave the counts unchanged.
As in the paper, pairs $(\pi,i)$ are counted once, not once for each
embedding they admit.

The decomposition argument differs from
Section~\ref{sec:full-master}.  Rather than construct the auxiliary
tree on the hyperedges, Lean proves directly that every tight tree,
with any chosen root edge, can be built from single edges by the
leaf and joining operations.  This is proved by induction from the
original edge order, including the steps needed to insert an edge
at any parent and to change the root.  The result is then used to
prove the main counting inequality.  The separate proof for paths
and the auxiliary-tree lemmas are not formalized line by line.

\subsection{Verification}

The project specifies the Lean version and includes the source files
and a script that compiles them and checks which axioms the theorems
depend on.  The theorem statements are collected in
\texttt{lean/Statements.lean}.  The permitted
logical axioms are \texttt{propext}, \texttt{Classical.choice}, and
\texttt{Quot.sound}.  Concrete formal tests cover root reorderings,
branches through different faces, empty hosts, and a sharp graph-path
example.

\begin{thebibliography}{99}

\bibitem{Kalai1984}
G.~Kalai,
Conjecture communicated in 1984, as recorded in
\cite[Conjecture~3.6 and reference~{[K2]}]{FF1987}.

\bibitem{FF1987}
P.~Frankl and Z.~F\"uredi,
\emph{Exact solution of some Tur\'an-type problems},
J. Combin. Theory Ser. A \textbf{45} (1987), 226--262.
\href{https://doi.org/10.1016/0097-3165(87)90016-1}
{doi:10.1016/0097-3165(87)90016-1}.

\bibitem{FJKMV2019}
Z.~F\"uredi, T.~Jiang, A.~Kostochka, D.~Mubayi, and J.~Verstra\"ete,
\emph{Hypergraphs not containing a tight tree with a bounded trunk},
SIAM J. Discrete Math. \textbf{33} (2019), 862--873.
\href{https://arxiv.org/abs/1712.04081}{arXiv:1712.04081}.

\bibitem{FJKMV2020}
Z.~F\"uredi, T.~Jiang, A.~Kostochka, D.~Mubayi, and J.~Verstra\"ete,
\emph{Tight paths in convex geometric hypergraphs},
Advances in Combinatorics (2020), Paper No.~1, 14 pp.
\href{https://doi.org/10.19086/aic.12044}{doi:10.19086/aic.12044}.
\href{https://arxiv.org/abs/2002.09457}{arXiv:2002.09457}, version~2.

\bibitem{BMP2025}
J.~P. Bright, K.~G. Milans, and J.~Porter,
\emph{Tur\'an numbers of ordered tight hyperpaths},
European J. Combin. \textbf{124} (2025), Article 104070.
\href{https://doi.org/10.1016/j.ejc.2024.104070}{doi:10.1016/j.ejc.2024.104070}.
\href{https://arxiv.org/abs/2212.13719}{arXiv:2212.13719}.

\bibitem{Stein2024}
M.~Stein,
\emph{Kalai's conjecture in $r$-partite $r$-graphs},
European J. Combin. \textbf{117} (2024), Article 103827.
\href{https://arxiv.org/abs/1912.11421}{arXiv:1912.11421}.

\bibitem{Erdos548Repository}
T.~Adamczewski (maintainer),
\emph{Erd\H{o}s problem \#548: Lean proof and provenance}, 2026.
\url{https://github.com/tadamcz/erdos548}.

\end{thebibliography}
\end{document}
